\section{Objetivo del Informe}

El objetivo central del presente informe es evaluar los resultados obtenidos tras la
implementación del sistema de gestión empresarial ERP Odoo 19 en Machu Picchu Exclusive
Tours E.I.R.L., analizando el impacto generado en sus procesos comerciales, operativos y de
seguridad de la información. La implementación tuvo como finalidad transformar el modelo
operativo tradicional de la empresa mediante la adopción de una plataforma integrada que
permita centralizar la información de los pasajeros, estandarizar el control de itinerarios
complejos y asegurar la trazabilidad del flujo de ventas.

Asimismo, se buscó dotar a la organización de una infraestructura tecnológica escalable que
fortalezca el control de cobros, reduzca los riesgos en el manejo de datos sensibles y mejore
la capacidad de toma de decisiones estratégicas basadas en evidencia. En este contexto, el
informe presenta una evaluación estructurada de los beneficios alcanzados, enfocándose en los
siguientes puntos clave:

\begin{itemize}
    \item \textbf{Centralización y Visión 360\textdegree:} evaluar la efectividad del CRM para unificar la base de datos de pasajeros (PAX) y eliminar la dispersión de información en chats de WhatsApp y hojas de Excel.
    \item \textbf{Seguridad y Privacidad de Datos:} verificar la correcta migración hacia un almacenamiento seguro de documentos de identidad (pasaportes) en el repositorio del sistema, mitigando los riesgos de ciberseguridad detectados en la fase de diagnóstico.
    \item \textbf{Control Operativo de Itinerarios:} analizar el desempeño del módulo de Proyectos y el tablero visual Kanban en la gestión de servicios que varían entre 2 a 20 días, asegurando un flujo ordenado desde la cotización hasta la ejecución final.
    \item \textbf{Trazabilidad Financiera:} medir la precisión en el registro manual de adelantos (del 30\% al 50\%) y saldos, garantizando el control interno necesario antes de proceder con la compra de boletos o servicios no reembolsables.
    \item \textbf{Nivel de Adopción y Eficiencia:} evidenciar el valor generado por la implementación en términos de reducción de errores administrativos, mejora en los tiempos de respuesta al cliente internacional y adopción de la herramienta por parte del equipo.
\end{itemize}
